\begin{theorem}[Curvature--regularity correspondence]
\label{thm:curv-reg-sc}
Let $K \subset B_R^D$ be a strictly convex body (boundary contains no line segment). Let $\Sigma \subset S^{D-1}$ be a closed set such that $\partial K$ is $C^2$ with positive Gauss curvature at all points $p^*(v)$ for $v \notin \Sigma$, and the maximum principal radius of curvature satisfies:
\[
  r_{\max}(v) \leq C_0 \cdot \mathrm{dist}(v, \Sigma)^{-\beta} \quad \text{for } v \notin \Sigma
\]
for some $\beta \in [0, 1)$. Then $h_K \in C^{1, 1-\beta}(S^{D-1})$.
\end{theorem}

\begin{proof}
The proof has three components.

\textbf{Component 1: The tangent-point map.}

For $v \in S^{D-1}$, define $p^*(v) = \arg\max_{x \in K} \langle x, v \rangle$. Since $K$ is strictly convex, the maximizer is unique for every $v$, and $p^* : S^{D-1} \to \partial K$ is a well-defined bijection (the reverse Gauss map). The gradient of $h_K$ exists everywhere and equals:
\[
  \nabla_{S^{D-1}} h_K(v) = p^*(v) - \langle p^*(v), v \rangle \, v = \Pi_{T_v S^{D-1}}(p^*(v)).
\]
This is a standard fact: $h_K(v) = \langle p^*(v), v \rangle$, and for nearby direction $v + \delta v$: $h_K(v + \delta v) \geq \langle p^*(v), v + \delta v \rangle = h_K(v) + \langle p^*(v), \delta v \rangle$, with equality to first order by optimality. Strict convexity of $K$ (unique maximizer) suffices.

\textbf{Component 2: Rate of change of the tangent-point map.}

For $v \notin \Sigma$ (where $\partial K$ is $C^2$ with positive curvature): the tangent-point map $p^*$ is differentiable, and by the implicit function theorem applied to the optimality condition $\nabla_{\partial K}(\langle \cdot, v \rangle) = 0$:
\[
  \|dp^*(v)\| = \|(\mathrm{II}_{p^*(v)})^{-1}\| = r_{\max}(v)
\]
where $\mathrm{II}$ is the second fundamental form. A small change $\delta v$ in the normal direction shifts the tangent point by $\delta p = (\mathrm{II})^{-1} \delta v$, and the eigenvalues of $(\mathrm{II})^{-1}$ are the principal radii $r_1, \ldots, r_{D-1}$.

Therefore, for $v \notin \Sigma$:
\[
  \|\nabla h_K(v_1) - \nabla h_K(v_0)\| \leq \|p^*(v_1) - p^*(v_0)\| + R |v_1 - v_0|
\]
and by the mean value theorem on the smooth part:
\[
  \|p^*(v_1) - p^*(v_0)\| \leq \sup_{t \in [0,1]} \|dp^*(\gamma(t))\| \cdot |v_1 - v_0| = \sup_t r_{\max}(\gamma(t)) \cdot |v_1 - v_0|
\]
along the geodesic $\gamma$ from $v_0$ to $v_1$.

\textbf{Component 3: Path integration and H\"older continuity.}

For \emph{any} $v_0, v_1 \in S^{D-1}$ (including directions in $\Sigma$), the gradient difference is bounded by integrating along the minimizing geodesic $\gamma : [0, d] \to S^{D-1}$:
\begin{align}
  \|\nabla h_K(v_1) - \nabla h_K(v_0)\| &\leq \int_0^d \left(r_{\max}(\gamma(t)) + R\right) dt \label{eq:path}
\end{align}
where the integrand is defined a.e.\ (it is defined for all $t$ with $\gamma(t) \notin \Sigma$, and $\Sigma$ has measure zero on $S^{D-1}$ since it is closed with empty interior by assumption).

Using $r_{\max}(v) \leq C_0 \cdot \mathrm{dist}(v, \Sigma)^{-\beta}$: the integral reduces to bounding $\int_0^d \mathrm{dist}(\gamma(t), \Sigma)^{-\beta} \, dt$.

\emph{Claim:} For any geodesic $\gamma$ of length $d$ on $S^{D-1}$ and any closed set $\Sigma$:
\[
  \int_0^d \mathrm{dist}(\gamma(t), \Sigma)^{-\beta} \, dt \leq \frac{C}{1-\beta} \cdot d^{1-\beta} \quad \text{for } \beta < 1.
\]

\emph{Proof of claim:} Let $\delta(t) = \mathrm{dist}(\gamma(t), \Sigma)$.

If $\gamma$ doesn't enter the $d$-neighborhood of $\Sigma$: $\delta(t) \geq \delta_{\min} > 0$ and the integral is $\leq d \cdot \delta_{\min}^{-\beta} < \infty$. Since $\delta_{\min} \geq $ const (the path stays far from $\Sigma$), this gives $O(d)$ which is $O(d^{1-\beta})$ for $d < 1$.

If $\gamma$ passes through the $d$-neighborhood of $\Sigma$: let $t_0$ be the point closest to $\Sigma$, with $\delta(t_0) = \delta_0$. By the triangle inequality on $S^{D-1}$: $\delta(t) \geq |\delta_0 - |t - t_0||$ (distance to $\Sigma$ changes by at most the path length). Therefore:
\begin{align*}
  \int_0^d \delta(t)^{-\beta} \, dt &\leq 2 \int_0^{d/2} (\delta_0 + s)^{-\beta} \, ds + 2\int_0^{d/2} s^{-\beta} \, ds \\
  &\leq 2 \int_0^{d/2} s^{-\beta} \, ds + 2\int_0^{d/2} s^{-\beta} \, ds \\
  &= 4 \cdot \frac{(d/2)^{1-\beta}}{1-\beta} = \frac{2^{2-\beta+1}}{1-\beta} \cdot \frac{d^{1-\beta}}{2} \leq \frac{C}{1-\beta} d^{1-\beta}.
\end{align*}
The worst case is $\delta_0 = 0$, where the path goes through $\Sigma$ itself; then $\delta(t) \geq |t - t_0|$ by the triangle inequality. \qed (claim)

Combining:
\[
  \|\nabla h_K(v_1) - \nabla h_K(v_0)\| \leq \frac{C \cdot C_0}{1-\beta} \cdot d(v_0, v_1)^{1-\beta} + R \cdot d(v_0, v_1).
\]
For $d(v_0, v_1) \leq 1$: the first term dominates (since $1-\beta < 1$), giving:
\[
  \|\nabla h_K(v_1) - \nabla h_K(v_0)\| \leq C' \cdot d(v_0, v_1)^{1-\beta}.
\]
Therefore $\nabla h_K \in C^{0, 1-\beta}(S^{D-1})$, hence $h_K \in C^{1, 1-\beta}(S^{D-1})$. \qed
\end{proof}

\begin{remark}[Extension to non-strictly-convex bodies ($\beta \geq 1$)]
For bodies with flat faces ($\beta \geq 1$): $K$ is not strictly convex, so $p^*(v)$ may not be unique at some directions, and $\nabla h_K$ does not exist in the classical sense at those directions. However:
\begin{itemize}
  \item $h_K$ is still Lipschitz on $S^{D-1}$ (with constant $R$).
  \item $\nabla h_K$ exists a.e.\ (Rademacher's theorem).
  \item At directions where $\nabla h_K$ does not exist, the subdifferential $\partial h_K(v)$ is a convex set (the face of $K$ with outward normal $v$).
  \item The ``jumps'' in $\nabla h_K$ at normal cone boundaries give $h_K \in C^{0,1}$ (Lipschitz) but not $C^{1,\alpha}$ for any $\alpha > 0$.
\end{itemize}
This corresponds to $\beta \geq 1 \Rightarrow h_K \in C^1$ at best, consistent with the $\beta$-rate $m^{-(2-\beta)/(D-1)} = m^{-1/(D-1)}$ for $\beta = 1$.
\end{remark}
